\documentclass[a4paper,10pt]{letter}
\usepackage[margin=0.8in]{geometry}
\usepackage[T1]{fontenc}
\usepackage{hyperref}
\hypersetup{
  hidelinks,
  pdftitle={Cover Letter: Fusion Utility Depends on the Evaluator},
  pdfauthor={Zigao Wu},
  pdfsubject={Submission to Infrared Physics \& Technology}
}
\setlength{\parskip}{0.45em}

\begin{document}
\sloppy
\begin{letter}{Editors\\Infrared Physics \& Technology}
\opening{Dear Editors,}

Please consider our manuscript, ``Fusion Utility Depends on the Evaluator:
A Boundary Study of Confidence-Shielded Input Selection for Infrared--Visible Object Detection,'' for
publication in \emph{Infrared Physics \& Technology}.

Infrared--visible fusion can expose complementary evidence, but the same fused
image may improve one detector and degrade another. We study an inference-time
policy that requires no ground-truth boxes. Confidence-Shielded,
Evaluator-Conditioned Fusion-Input Selection (CS-ECF) scores visible input,
consensus-detail fusion (CDF), PIAFusion, and SPGFusion with the target
detector. It uses a fused input only when its top-five confidence reaches 1.5
times visible; otherwise, visible is retained.

The contribution is a bounded empirical result, not a claim of general
accuracy improvement. On the scene-disjoint LLVIP evaluation, the primary
visible contrast was \(+0.00012\) \([-0.00484,+0.00670]\) and CS-ECF did not
outperform Always CDF. Relative to the unconstrained no-shield selector, its
mean AP50 gain was \(+0.00894\) \([+0.00096,+0.02155]\). A random allocation
matched to the same switch count and candidate mix was worse. On VEDAI,
both training runs favoured the shield over no shield, with an equal-run gain
of \(+0.0410\) \([+0.0097,+0.0782]\), while CS-ECF remained below visible
input by \(-0.0136\). The apparent safeguard was specific to AP50 and did not
persist under AP50--95.

We also trained the public YOLOv11-RGBT early-fusion detector on the same M3FD
split as a single-seed implementation pilot. It provides a recent end-to-end
system context but not a tuned performance ranking. External evaluators
reversed both the visible-relative and no-shield orderings. Matched
image-statistic controls did not establish a confidence-specific mechanism.

The manuscript therefore supports CS-ECF only as an AP50 safeguard after
detector-specific validation. It recommends visible fallback or a stable CDF
as the default and a switch-budget-matched random check before deployment.
The new case-study figure shows one successful and one failed switch selected
by a fixed rule.

The manuscript is original, has not been published previously, and is not under
consideration by another journal. It is a substantially revised successor to
manuscript INFPHY-D-26-02136, which the same authors previously submitted to
this journal and which was not accepted. The earlier study examined
evaluator-dependent rankings of fidelity metrics. The present study instead
evaluates an inference-time selection policy and does not reproduce the earlier
fidelity-ranking claims. To address the principal concerns raised in prior
review, we added scene-disjoint LLVIP and VEDAI evaluation, independent training
runs, standard confidence-ordered AP scoring, a switch-budget and candidate-mix
matched random control, external-evaluator stress tests, a recent joint
fusion-detection baseline, explicit operational conditions, and complete
supplementary evidence. We also narrowed every claim to the observed AP50
setting and report the null visible contrast, negative CDF comparison, and
external reversal directly. All shared datasets and evaluator infrastructure
are identified in the manuscript. All authors have approved the submission and
declare no competing interests. This research did not receive any specific
grant from funding agencies in the public, commercial, or not-for-profit
sectors. The manuscript includes an AI-use declaration and statements for data,
code, CRediT contributions, and competing interests.

Thank you for your consideration.

\closing{Sincerely,}
Zigao Wu\\
North China Electric Power University\\
Baoding, Hebei 071003, China\\
\href{mailto:zgwu@ncepu.edu.cn}{zgwu@ncepu.edu.cn}
\end{letter}
\end{document}
